\documentclass[
    language=english,
    doctype=clinical-report,
    institution=none,
]{../../omnilatex}

\RequirePackage{config/document-settings}

\begin{document}
\maketitle

\begin{abstract}
\textbf{Title:} A Phase III, Randomized, Double-Blind, Placebo-Controlled Trial
of Mirabegron 50 mg Once Daily for the Treatment of Overactive Bladder Syndrome
in Adult Patients: 12-Week Efficacy and Safety Results.

\textbf{Protocol ID:} MIRA-OAB-301

\textbf{Sponsor:} PharmaCure Therapeutics, Inc.

\textbf{ClinicalTrials.gov Identifier:} NCT05892341

\textbf{Objective:} To evaluate the efficacy and safety of mirabegron 50 mg
once daily compared to placebo in adult patients with symptoms of overactive
bladder (OAB) inadequately managed with prior antimuscarinic therapy.

\textbf{Methods:} Adult patients ($\geq$ 18 years) with OAB symptoms for
$\geq$ 3 months were randomized 1:1 to mirabegron 50 mg ($n = 462$) or
placebo ($n = 458$) once daily for 12 weeks. Primary endpoints were change
from baseline in number of micturitions per 24 hours and number of urgency
episodes per 24 hours at week 12.

\textbf{Results:} Mirabegron significantly reduced micturitions per 24 hours
versus placebo ($-1.87$ vs.\ $-1.18$; difference $-0.69$; 95\% CI $[-0.89,
-0.49]$; $p < .001$) and urgency episodes ($-2.31$ vs.\ $-1.42$; difference
$-0.89$; 95\% CI $[-1.14, -0.64]$; $p < .001$). Treatment-emergent adverse
events occurred in 34.2\% of mirabegron-treated patients and 28.4\% of
placebo-treated patients. The most common adverse events were urinary tract
infection (5.6\%), nasopharyngitis (4.3\%), and headache (3.2\%).

\textbf{Conclusions:} Mirabegron 50 mg once daily demonstrated statistically
significant and clinically meaningful improvements in OAB symptoms versus
placebo, with a safety profile consistent with the established mirabegron
label.
\end{abstract}

\section{Introduction}
\label{sec:introduction}

Overactive bladder (OAB) is a chronic condition characterized by urgency,
with or without urge incontinence, typically accompanied by frequency and
nocturia \cite{haylen2010international}. OAB affects approximately 12--16\%
of the adult population and significantly impacts quality of life, work
productivity, and healthcare utilization \cite{irwin2011epidemiology}.

Mirabegron is a selective $\beta_3$-adrenergic receptor agonist approved
for the treatment of OAB symptoms. Unlike antimuscarinic agents, mirabegron
works by relaxing the detrusor muscle during the storage phase, offering an
alternative mechanism of action with a different side effect profile
\cite{chapple2014randomised}.

The purpose of this Phase III trial was to confirm the efficacy and safety
of mirabegron 50 mg once daily in a large, well-characterized OAB population,
including patients who had inadequate response to or experienced intolerable
side effects with prior antimuscarinic therapy.

\section{Methods}
\label{sec:methods}

\subsection{Study Design}

This was a 12-week, randomized, double-blind, placebo-controlled, multicenter
trial conducted at 98 sites across the United States, Europe, and Asia-Pacific
between March 2023 and November 2024.

\subsection{Patient Population}

Eligible patients were adults $\geq$ 18 years of age with symptoms of OAB
(urgency with or without incontinence, frequency $\geq$ 8 micturitions/24 h,
and $\geq$ 3 urgency episodes/24 h) for $\geq$ 3 months. Patients must have
failed $\geq$ 1 prior antimuscarinic therapy or discontinued due to
intolerable side effects (e.g., dry mouth, constipation, blurred vision).

\begin{table}[htbp]
    \centering
    \caption{Key inclusion and exclusion criteria.}
    \label{tab:criteria}
    \begin{tabular}{@{}p{7cm}p{7cm}@{}}
        \toprule
        \textbf{Inclusion Criteria} & \textbf{Exclusion Criteria} \\
        \midrule
        Age $\geq$ 18 years & Neurogenic OAB or detrusor overactivity \\
        OAB symptoms $\geq$ 3 months & Urinary retention (post-void residual $>$ 200 mL) \\
        Mean micturitions $\geq$ 8/24 h on voiding diary & Active urinary tract infection \\
        Mean urgency episodes $\geq$ 3/24 h & History of bladder cancer or radiation therapy \\
        Failed $\geq$ 1 prior antimuscarinic & Uncontrolled hypertension ($>$ 160/100 mmHg) \\
        Able to complete 3-day voiding diary & Pregnant, nursing, or planning pregnancy \\
        \bottomrule
    \end{tabular}
\end{table}

\subsection{Randomization and Blinding}

Patients were randomized 1:1 to mirabegron 50 mg or matching placebo once
daily, stratified by:
\begin{itemize}
    \item Baseline micturition frequency ($<$ 10 vs.\ $\geq$ 10 per 24 h)
    \item Prior antimuscarinic use (yes vs.\ no)
    \item Geographic region (North America vs.\ Europe vs.\ Asia-Pacific)
\end{itemize}

The randomization sequence was generated using a validated interactive
response technology (IRT) system. Blinding was maintained through identical
tablet appearance and packaging.

\subsection{Endpoints}

\subsubsection{Primary Endpoints}
\begin{enumerate}
    \item Change from baseline to week 12 in mean number of micturitions
          per 24 hours (from 3-day voiding diary).
    \item Change from baseline to week 12 in mean number of urgency episodes
          per 24 hours (from 3-day voiding diary).
\end{enumerate}

\subsubsection{Key Secondary Endpoints}
\begin{itemize}
    \item Change from baseline to week 12 in mean volume voided per micturition.
    \item Change from baseline to week 12 in number of incontinence episodes
          per 24 hours (in patients with incontinence at baseline).
    \item Proportion of patients with $\geq$ 50\% reduction in incontinence
          episodes.
    \item Patient Perception of Bladder Condition (PPBC) score change.
    \item Treatment satisfaction (TS-POSS) score change.
\end{itemize}

\subsection{Safety Assessments}

Safety was assessed by monitoring treatment-emergent adverse events (TEAEs),
vital signs, 12-lead ECG, clinical laboratory parameters, and post-void
residual volume.

\section{Results}
\label{sec:results}

\subsection{Patient Disposition}

\begin{figure}[htbp]
    \centering
    \begin{tikzpicture}[
        node distance=1.2cm,
        box/.style={rectangle, draw, text width=4.5cm, align=center, minimum height=0.9cm},
        arrow/.style={->, thick}
    ]
        \node[box] (screened) {Patients screened ($n = 1,247$)};
        \node[box, below of=screened] (randomized) {Patients randomized ($n = 920$)};
        \node[box, below left=1.2cm and 0.5cm of randomized] (mirab) {Mirabegron 50 mg ($n = 462$)};
        \node[box, below right=1.2cm and 0.5cm of randomized] (placebo) {Placebo ($n = 458$)};
        \node[box, below of=mirab] (mirab-completed) {Completed 12 weeks ($n = 438$)};
        \node[box, below of=placebo] (placebo-completed) {Completed 12 weeks ($n = 441$)};

        \draw[arrow] (screened) -- (randomized);
        \draw[arrow] (randomized) -- node[left, sloped, above] {1:1} (mirab);
        \draw[arrow] (randomized) -- node[right, sloped, above] {1:1} (placebo);
        \draw[arrow] (mirab) -- node[left] {24 discontinued} (mirab-completed);
        \draw[arrow] (placebo) -- node[right] {17 discontinued} (placebo-completed);
    \end{tikzpicture}
    \caption{Patient disposition flow diagram.}
    \label{fig:disposition}
\end{figure}

\subsection{Baseline Characteristics}

\begin{table}[htbp]
    \centering
    \caption{Baseline demographic and clinical characteristics.}
    \label{tab:baseline}
    \begin{tabular}{@{}lcc@{}}
        \toprule
        Characteristic & Mirabegron 50 mg ($n = 462$) & Placebo ($n = 458$) \\
        \midrule
        Age, mean $\pm$ SD (years)         & 58.3 $\pm$ 12.4 & 57.9 $\pm$ 12.1 \\
        Female, $n$ (\%)                    & 371 (80.3)       & 364 (79.5) \\
        BMI, mean $\pm$ SD (kg/m$^2$)      & 29.4 $\pm$ 5.8   & 29.1 $\pm$ 5.6 \\
        Micturitions/24 h, mean $\pm$ SD   & 10.8 $\pm$ 2.3   & 10.9 $\pm$ 2.4 \\
        Urgency episodes/24 h, mean $\pm$ SD & 5.4 $\pm$ 2.1 & 5.3 $\pm$ 2.0 \\
        Incontinence episodes/24 h, mean $\pm$ SD & 3.2 $\pm$ 1.8 & 3.1 $\pm$ 1.7 \\
        Duration of OAB, median (IQR) (years) & 4.2 (2.1--7.8) & 4.0 (2.0--7.5) \\
        Prior antimuscarinic use, $n$ (\%)  & 289 (62.6)       & 278 (60.7) \\
        \bottomrule
    \end{tabular}
\end{table}

\subsection{Primary Efficacy Endpoints}

\begin{table}[htbp]
    \centering
    \caption{Primary efficacy results at week 12 (ITT population).}
    \label{tab:primary-efficacy}
    \begin{tabular}{@{}lcccc@{}}
        \toprule
        Endpoint & Treatment & Baseline & Week 12 & LS Mean Change \\
        \midrule
        \multirow{2}{*}{Micturitions/24 h} & Mirabegron & 10.8 & 8.9 & $-1.87$ \\
        & Placebo & 10.9 & 9.7 & $-1.18$ \\[6pt]
        \multirow{2}{*}{Urgency episodes/24 h} & Mirabegron & 5.4 & 3.1 & $-2.31$ \\
        & Placebo & 5.3 & 3.9 & $-1.42$ \\[6pt]
        \multirow{2}{*}{Volume voided (mL)} & Mirabegron & 172 & 214 & $+42$ \\
        & Placebo & 174 & 196 & $+22$ \\
        \midrule
        \multicolumn{5}{l}{\textbf{Treatment differences (Mirabegron $-$ Placebo)}} \\
        \midrule
        Micturitions/24 h        & & & & $-0.69$ \\
        Urgency episodes/24 h    & & & & $-0.89$ \\
        Volume voided (mL)       & & & & $+20$ \\
        \bottomrule
    \end{tabular}
\end{table}

Both primary endpoints met statistical significance ($p < .001$). The
treatment difference for micturitions was $-0.69$ per 24 hours (95\% CI
$[-0.89, -0.49]$) and for urgency episodes was $-0.89$ per 24 hours
(95\% CI $[-1.14, -0.64]$).

\subsection{Key Secondary Endpoints}

In patients with incontinence at baseline ($n = 523$), mirabegron reduced
incontinence episodes by 2.4 per 24 hours compared to 1.5 with placebo
(treatment difference $-0.9$; $p < .001$). The proportion of patients
achieving $\geq$ 50\% reduction in incontinence was 68.2\% for mirabegron
versus 49.1\% for placebo ($p < .001$).

\subsection{Safety Results}

\begin{table}[htbp]
    \centering
    \caption{Treatment-emergent adverse events occurring in $\geq$ 2\% of
             patients in either group.}
    \label{tab:safety}
    \begin{tabular}{@{}lcc@{}}
        \toprule
        Adverse Event & Mirabegron 50 mg ($n = 462$) & Placebo ($n = 458$) \\
        \midrule
        Any TEAE & 158 (34.2) & 130 (28.4) \\
        Urinary tract infection & 26 (5.6) & 22 (4.8) \\
        Nasopharyngitis & 20 (4.3) & 15 (3.3) \\
        Headache & 15 (3.2) & 11 (2.4) \\
        Dry mouth & 12 (2.6) & 3 (0.7) \\
        Constipation & 10 (2.2) & 8 (1.7) \\
        Diarrhea & 9 (1.9) & 11 (2.4) \\
        Nausea & 8 (1.7) & 6 (1.3) \\
        Fatigue & 7 (1.5) & 5 (1.1) \\
        \midrule
        Serious TEAEs & 8 (1.7) & 6 (1.3) \\
        TEAEs leading to discontinuation & 14 (3.0) & 7 (1.5) \\
        \bottomrule
    \end{tabular}
\end{table}

Blood pressure increases of $>$ 15/10 mmHg were observed in 8.4\% of
mirabegron-treated patients versus 4.1\% of placebo-treated patients.
Mean change in systolic blood pressure from baseline was $+2.1$ mmHg
(mirabegron) versus $+0.4$ mmHg (placebo).

\section{Discussion}
\label{sec:discussion}

The results of this Phase III trial demonstrate that mirabegron 50 mg
once daily provides statistically significant and clinically meaningful
improvements in OAB symptoms compared to placebo. The treatment effects
observed for both primary endpoints ($-0.69$ micturitions and $-0.89$
urgency episodes) exceed the minimally important difference thresholds
established in the literature ($-0.5$ for micturitions and $-1.0$ for
urgency episodes) \cite{noble2016defining}.

The safety profile was consistent with the established mirabegron label.
The slightly higher incidence of dry mouth compared to placebo (2.6\%
vs.\ 0.7\%) confirms the $\beta_3$-agonist mechanism avoids the
anticholinergic side effects that frequently lead to discontinuation
of antimuscarinic agents.

\section{Conclusions}
\label{sec:conclusions}

Mirabegron 50 mg once daily is an effective and well-tolerated treatment
for OAB in adults, including those who have failed prior antimuscarinic
therapy. These results support mirabegron as a first-line pharmacological
option for OAB management.

\bibliographystyle{apa}
\bibliography{references}

\end{document}
